\chapter{MLOps Implementation}
\label{ch:mlops}

\section{System Architecture}

Figure~\ref{fig:architecture} shows the end-to-end CAAS architecture.
Data flows from three external sources through an automated ingestion pipeline
into a feature store (Amazon S3), is used to train and register models in
MLflow, served via FastAPI on EC2, and monitored continuously by Evidently.

\begin{figure}[H]
    \centering
    \includegraphics[width=\textwidth]{fig_cloud_architecture}
    \caption{CAAS cloud architecture on AWS. Three external APIs (air4thai,
    Open-Meteo, NASA FIRMS) feed a GitHub Actions-orchestrated ingestion layer
    (3-hourly cron) into four S3 tiers (Raw, Processed, Feature Store, Model
    Artifacts). A single EC2~t3.micro instance hosts FastAPI, MLflow, and
    Streamlit. Evidently AI drift monitoring runs as a daily scheduled job.}
    \label{fig:architecture}
\end{figure}

\section{Experiment Tracking and Model Registry}

All training runs are logged to MLflow with:
\begin{itemize}
    \item Full hyperparameter sets per horizon.
    \item Validation and test metrics (MAE, RMSE, R², alert F1, AUROC).
    \item Model artefacts (\texttt{.json} for XGBoost, \texttt{.keras} for LSTM).
    \item Feature importance and prediction CSVs as run artefacts.
\end{itemize}

Three experiments are maintained: \texttt{CAAS-XGBoost} (3 runs),
\texttt{CAAS-LSTM} (3 runs), and \texttt{CAAS-Ablation-NoFIRMS} (3 runs).

\begin{figure}[H]
    \centering
    \includegraphics[width=0.85\textwidth]{fig_mlflow_runs}
    \caption{MLflow experiment view showing CAAS-XGBoost runs.}
    \label{fig:mlflow_runs}
\end{figure}

\section{Model Serving}

The champion XGBoost model is served via a FastAPI application
(\texttt{04\_Scripts/serve/app.py}) exposing five endpoints:

\begin{table}[H]
\centering
\caption{CAAS API Endpoints}
\begin{tabular}{llp{7cm}}
\toprule
\textbf{Method} & \textbf{Endpoint} & \textbf{Description} \\
\midrule
GET  & \texttt{/health}      & Liveness check, lists loaded models \\
GET  & \texttt{/forecast}    & Latest t+1/3/7 forecasts from most recent feature row \\
POST & \texttt{/predict}     & On-demand prediction from supplied feature values \\
GET  & \texttt{/history}     & Recent PM2.5 history for dashboard charting \\
GET  & \texttt{/model/info}  & Champion model metadata and test metrics \\
\bottomrule
\end{tabular}
\end{table}

The service is containerised using a multi-stage Docker build that reduces
the image size to approximately 300~MB by excluding training and monitoring
dependencies. On EC2, the container starts automatically via a user-data
bootstrap script that pulls the latest model artefacts from S3 at launch.

A Streamlit dashboard (\texttt{dashboard.py}) provides a public-facing
interface showing colour-coded forecast cards, a 60-day PM2.5 history chart,
and the model performance summary.

\begin{figure}[H]
    \centering
    \includegraphics[width=0.85\textwidth]{fig_streamlit_dashboard}
    \caption{CAAS Streamlit dashboard showing PM2.5 forecast cards and history.}
    \label{fig:streamlit}
\end{figure}

\section{CI/CD Pipeline}

Three GitHub Actions workflows implement the CI/CD pipeline:

\textbf{test.yml} runs on every push to \texttt{main}:
\begin{enumerate}
    \item Install dependencies.
    \item Lint with \texttt{ruff}.
    \item Run 38 pytest tests covering feature schema, model inference, and API endpoints.
\end{enumerate}

\textbf{daily\_pipeline.yml} runs on cron (06:00 Bangkok time):
\begin{enumerate}
    \item Fetch latest PM2.5, weather, and FIRMS data.
    \item Rebuild feature matrix.
    \item Run inference with champion model.
    \item Upload results to S3.
    \item Run Evidently drift check; trigger retraining if drift thresholds are exceeded.
\end{enumerate}

\textbf{retrain.yml} triggered by drift alert or annual schedule (December~1):
\begin{enumerate}
    \item Full data snapshot and feature rebuild.
    \item Train new XGBoost candidate models (all three horizons).
    \item \textbf{Validation gate:} candidate must achieve $\geq$5\% MAE improvement
          and alert F1 $\geq$ 0.75 over a 90-day rolling window.
    \item If gate passes: promote via MLflow (\texttt{Staging → Production}).
    \item Restart FastAPI on EC2 to serve the new champion.
\end{enumerate}

\begin{figure}[H]
    \centering
    \includegraphics[width=0.85\textwidth]{fig_model_promotion}
    \caption{Champion/challenger model promotion flow with validation gate.}
    \label{fig:model_promotion}
\end{figure}

\section{Monitoring and Drift Detection}

The Evidently-based drift monitor (\texttt{monitoring/evidently\_report.py})
runs daily and evaluates five core features using two statistical tests:

\begin{itemize}
    \item \textbf{PSI (Population Stability Index):} Detects distributional shift.
          Threshold: PSI $>$ 0.20.
    \item \textbf{KS test:} Two-sample Kolmogorov-Smirnov test.
          Threshold: $p < 0.05$.
\end{itemize}

The retrain trigger policy (updated 2026-04-03) requires:
\begin{itemize}
    \item Immediate retrain if 7-day rolling MAE exceeds 15~\textmu g/m\textsuperscript{3}, \textbf{or}
    \item Both core features (\texttt{pm25\_lag1}, \texttt{pm25\_roll7\_mean}) show drift, \textbf{or}
    \item One core feature and two or more soft features show drift.
\end{itemize}

Seasonal KS references are used for \texttt{wind\_speed}, \texttt{hotspot\_50km},
and \texttt{is\_haze\_season} to avoid spurious drift flags from expected
seasonal variation. Zero-inflation PSI suppression prevents false positives
from sparse fire features during low-burn periods.

\begin{figure}[H]
    \centering
    \includegraphics[width=0.85\textwidth]{fig_drift_output}
    \caption{Evidently drift report output (2026-04-03). Core drift detected
    in \texttt{pm25\_lag1} and \texttt{pm25\_roll7\_mean} during December 2025 ---
    expected seasonal behaviour, not model failure.}
    \label{fig:drift}
\end{figure}

When drift is detected with \texttt{--trigger-retrain}, the monitor fires a
GitHub \texttt{repository\_dispatch} event (\texttt{drift-alert}), which
automatically triggers \texttt{retrain.yml} to begin the retraining pipeline.

\section{Infrastructure as Code}

All AWS infrastructure is defined in Terraform (\texttt{infra/}):

\begin{itemize}
    \item \textbf{S3 bucket} with versioning, server-side encryption (AES-256),
          and public access block.
    \item \textbf{IAM role and policy} granting EC2 read/write access to S3
          without embedding credentials in code.
    \item \textbf{Security group} opening ports 8000 (API), 5000 (MLflow), and 22 (SSH).
    \item \textbf{EC2 t3.micro instance} with user-data bootstrap: installs Docker,
          pulls the model artefacts from S3, and starts the FastAPI container.
          FastAPI, Streamlit, and MLflow are co-located on this single instance.
\end{itemize}

Provisioning requires only \texttt{terraform init \&\& terraform apply}.
\texttt{terraform destroy} cleanly removes all resources when not in use,
avoiding unnecessary cost.
